\documentclass{article}
\usepackage{iclr2025_conference,times}

\input{math_commands.tex}

\usepackage{amsmath}
\usepackage{booktabs}
\usepackage{graphicx}
\usepackage{hyperref}
\usepackage{microtype}
\usepackage{url}
\usepackage{xcolor}
\usepackage{adjustbox}
\usepackage{multirow}
\usepackage{booktabs}
\usepackage{caption}

\title{Llama Scope 2: Sparse Dictionaries, Replacement Models,\\
and the Llamascopium Framework}

\author{Llama Scope Team\\
OpenMOSS\\
\texttt{TODO: contact email}}

\newcommand{\todo}[1]{\textcolor{red}{[TODO: #1]}}

\iclrfinalcopy

\begin{document}

\maketitle

\begin{abstract}
Sparse autoencoders (SAEs) and their variants provide sparse dictionaries for 
interpreting language model activations in mechanistic interpretability. Furthermore, 
Transcoders and Low-Rank Sparse Attention (Lorsa) modules provide methods for 
decomposing MLP and attention computations into interpretable sparse dictionaries.
We introduce Llama Scope 2, a unified suite of per-layer
Transcoders and Lorsas modules for the Qwen3,
Llama 3.2, and Gemma 3 model families. This release fills a practical gap in
current circuit-analysis infrastructure, where popular open language models
lack complete and consistent sparse dictionaries for both MLP and attention
computations. We evaluate these dictionaries in terms of reconstruction
quality, faithfulness, generalization, and scalability, and use case studies to
show how they can be composed into replacement models for complete circuit
analysis. The checkpoints for the Llama Scope 2 Sparse Dictionary Families are
publicly available at \todo{checkpoint URL}, and the training, analysis, and
visualization framework is released at \todo{code URL}. Together, these
resources provide an open foundation for broad circuit-level analysis and model
biology studies of modern language models.
\end{abstract}

\section{Introduction}
\label{sec:introduction}

Sparse autoencoders (SAEs) learn sparse and overcomplete dictionaries in which a model
activation is represented by a small number of active features \citep{bricken2023monosemanticity, cunningham2023sae}. 
These features provide a potentially interpretable decomposition of the information represented within language models 
and have become foundational tools for research on language model interpretability 
\citep{he2024llamascope, lieberum2024gemma, deng2026qwenscope}.

Furthermore, Transcoders decompose MLP computations, enabling preliminary circuit analysis through 
MLP layers \citep{dunefsky2024transcoder}. 
Low-Rank Sparse Attention (Lorsa) modules extend this sparse decomposition to attention modules, 
providing interpretable sparse dictionaries for token-level interactions within attention \citep{he2025lorsa}.
Together, these dictionaries enable a sparse decomposition , 
and support circuit  analysis across complete Transformers and provide an interpretable 
view of token-level feature interactions \citep{shu2026completereplacement}.

The original Llama Scope release scaled this approach to every
layer and sublayer of Llama-3.1-8B and provided a large collection of
ready-to-use TopK SAEs \citep{he2024llamascope}. This reduced the cost of
retraining dictionaries and made it easier to compare sparse features across
layers, sites, and dictionary widths.

Llama Scope 2 brings these components together in a unified release spanning Qwen3, Llama 3.2, and 
Gemma 3 \citep{deng2026qwenscope, mcdougall2025gemmascope2}. We use TopK training for the Qwen3 and Llama 3.2 families, 
and JumpReLU training for Gemma 3 Lorsas following the recipe adopted by Gemma Scope 2 \citep{mcdougall2025gemmascope2}. 
The release is accompanied by the new version of Llamascopium, a distributed framework that unifies the complete sparse 
interpretability workflow, including activation generation, sparse dictionary training, feature analysis, visualization, and circuit analysis.

To make circuit analysis more flexible and customizable, we further train
Matryoshka SAEs \citep{bussmann2025matryoshka} on the last residual stream in the circuit-tracing pipeline.
The resulting hierarchical features allow users to select a residual-feature range of interest while suppressing 
irrelevant graph components, providing a way to address \emph{Uninterested Circuits} \citep{shu2026completereplacement}.

In section~\ref{sec:background}, we introduce the components of Llama Scope 2: Transcoders, and Low-Rank Sparse Attention (Lorsa).
% In section~\ref{sec:llamascopium}, we introduce the Llamascopium framework, which provides a unified interface for training, analyzing, and visualizing the sparse dictionaries.
% In section~\ref{sec:release}, we introduce the Llama Scope 2 release, which includes the sparse dictionaries for Qwen3, Llama 3.2, and Gemma 3.
In section~\ref{sec:training}, we introduce the training procedure for the sparse dictionaries.
In section~\ref{sec:evaluation}, we evaluate the performance of the sparse dictionaries.
In section~\ref{sec:attribution_graph_construction}, we construct the attribution graph for the sparse dictionaries.
In section~\ref{sec:case_studies}, we demonstrate the use of the sparse dictionaries for case studies.
In section~\ref{sec:attribution_graph_evaluation}, we evaluate the quality of the attribution graph.


\paragraph{Contributions.}
This technical report documents the following components of Llama Scope 2:
\begin{itemize}
    \item \textbf{A multi-family sparse dictionary suite.}
    We train layer-wise transcoders and Lorsas for Llama 3.2 1B and 3B, Qwen3 1.7B and
    8B, and Gemma 3 270M, 1B, and 4B.

    \item \textbf{Replacement models for circuit analysis.}
    These components can be composed into sparse replacement models that expose
    feature-to-feature computational pathways while retaining explicit
    reconstruction-error terms.

    \item \textbf{Matryoshka SAE for uninterested circuits.}
    We train Matryoshka SAEs on the last residual stream to define hierarchical
    residual features, enabling \emph{Uninterested Circuits} that prune circuit
    graphs by selecting a desired feature range.

    \item \textbf{The Llamascopium framework.}
    We provide a unified, distributed codebase for activation generation,
    sparse-dictionary training, feature analysis, visualization, and circuit
    attribution.
\end{itemize}

\section{Sparse Dictionary Background}
\label{sec:background}

\subsection{Vanilla Sparse Autoencoders}
\label{sec:vanilla_sae}

Given a model activation $\vx \in \mathbb{R}^{d_{\mathrm{model}}}$, a sparse
autoencoder encodes and reconstructs it as
\begin{equation}
\begin{aligned}
    \vh &= \mW_E \vx + \vb_E, \qquad
    \vz = \sigma(\vh), \\
    \widehat{\vx} &= \mW_D \vz + \vb_D.
\end{aligned}
\end{equation}
Here, $\vh$ denotes the encoder pre-activations and $\vz$ the sparse feature
activations. The columns of $\mW_D$ form an overcomplete dictionary of feature
directions, while the nonzero entries of $\vz$ specify how these directions are
linearly combined to reconstruct $\vx$ (up to the decoder bias $\vb_D$).
Llama Scope 2 primarily uses two activation functions: TopK, which directly controls the 
number of active latents per token, and JumpReLU, which learns a separate activation threshold 
for each latent \citep{gao2024topk, rajamanoharan2024jumprelu}.

\subsection{Per-layer Transcoders}
\label{sec:per_layer_transcoder}

Per-layer transcoders have the same encoder--decoder structure as sparse autoencoders, 
but differ in their training objective \citep{dunefsky2024transcoder}. Rather than reconstructing a representation from itself, 
a transcoder takes the normalized input to an MLP sublayer and predicts the corresponding MLP output. 
In this sense, transcoders provide a sparse decomposition of the MLP computation, 
rather than of a single representation.

For each layer $\ell$, the transcoder maps the normalized MLP input
$\vx_{\mathrm{MLP,in}}$ to the corresponding MLP output
$\vx_{\mathrm{MLP,out}}$ as
\begin{equation}
    \begin{gathered}
        \vh = \mW_E \vx_{\mathrm{MLP,in}} + \vb_E,
        \qquad
        \vz = \sigma(\vh), \\
        \widehat{\vx}_{\mathrm{MLP,out}}
        = \mW_D \vz + \vb_D.
    \end{gathered}
\end{equation}


\subsection{Low-rank Sparse Attention (Lorsa)}
\label{sec:lorsa}

Lorsa can be viewed as a transcoder-style decomposition of attention layers \citep{he2025lorsa}. 
It represents attention as a sparse dictionary of rank-1 features, where the OV 
weights serve as the decoder and decoder, while the QK pattern determines feature activation and 
represents the interactions across token positions.

For an attention input
$\vx_{\mathrm{Attn,in}} \in \mathbb{R}^{n \times d_{\mathrm{model}}}$,
each Lorsa feature $k$ first computes its attention pattern as
\begin{equation}
\begin{gathered}
    \mA_k =
    \operatorname{softmax}\!\left(
    \frac{
        (\vx_{\mathrm{Attn,in}}\mW_{Q,k})
        (\vx_{\mathrm{Attn,in}}\mW_{K,k})^\top
    }{\sqrt{d_h}}
    + \mB
    \right), \\[4pt]
    h_{i,k}
    =
    \sum_{j=1}^{n}
    A_{k,i,j}\,
    \vw_{V,k}^{\top}
    \vx_{\mathrm{Attn,in},j}.
\end{gathered}
\end{equation}

The sparse feature activation is obtained by applying the activation function
\begin{equation}
    z_{i,k} = \sigma(h_{i,k}).
\end{equation}

Finally, the active features are decoded into the attention output as
\begin{equation}
    \widehat{\vx}_{\mathrm{Attn,out},i}
    =
    \sum_k z_{i,k}\vw_{O,k}
    + \vb_O.
\end{equation}
Here, $\vz_i \in \mathbb{R}^{n \times 1}$ denotes the scalar activation of feature $i$ across token positions, 
and $\vw_{O,i}$ specifies its output direction. In this sense, Lorsa provides a sparse decomposition of the 
attention computation.

\subsection{Matryoshka Sparse Autoencoders}
\label{sec:matryoshka_sae}

Matryoshka Sparse Autoencoders (Matryoshka SAEs) extend standard SAEs by
organizing the latent space into a sequence of nested prefixes
\citep{bussmann2025matryoshka}. Each prefix is trained to reconstruct the same
activation, encouraging earlier features to capture more general structure and
later features to provide progressively finer-grained representations.

Given an activation $\vx$, the encoder is the same as in a standard SAE,
\begin{equation}
\begin{gathered}
    \vh = \mW_E \vx + \vb_E,
    \qquad
    \vz = \sigma(\vh),
\end{gathered}
\end{equation}
while different prefixes $\vz^{(m)}$ of the latent vector are decoded independently,
\begin{equation}
    \widehat{\vx}^{(m)}
    =
    \mW_D \vz^{(m)} + \vb_D,
    \qquad m \in \mathcal{M},
\end{equation}

where $\mathcal{M}$ denotes the set of predefined prefix widths used to form the nested latent representations.
We train Matryoshka SAEs on the final residual stream of the model, 
immediately before the unembedding layer. We find that they can organize concepts into hierarchical levels, 
for example, from ``person name'' to ``female name'' and finally to a specific name such as ``Mary.'' 
This hierarchy allows circuit tracing to selectively identify circuits at different levels of abstraction.

\section{Training Llama Scope 2}
\label{sec:training}

\subsection{Release Overview}

Table~\ref{tab:release} summarizes the planned organization of the release.
The table records the configurations currently represented by the main
training workflows. The final public-checkpoint count and any additional SAE
or Matryoshka variants will be added after the final checkpoint audit.

\begin{table*}[t]
    \caption{
        Core Llama Scope 2 training configurations. TopK values are per token.
        For JumpReLU Lorsas, $L_0^\star$ is layer-dependent; arrows denote the
        target-$L_0$ schedules from early to deeper layers.
        Each listed configuration is trained layer-wise.
        }
    \centering
    \small
    \resizebox{\textwidth}{!}{%
    \begin{tabular}{lrrllll}
        \toprule
        Backbone & Layers & $d_{\mathrm{model}}$ & Dictionary & Width &
        Activation / sparsity & Training tokens \\
        \midrule

        Llama 3.2 1B & 16 & 2048 & TC, Lorsa & $8\times$, $32\times$ &
        TopK, $K\in\{64,128,256\}$ & 8B \\

        Llama 3.2 3B & 28 & 3072 & TC, Lorsa & $8\times$, $32\times$ &
        TopK, $K\in\{64,128,256\}$ & 8B \\

        \midrule

        Qwen3 1.7B & 28 & 2048 & TC, Lorsa & $8\times$, $32\times$ &
        TopK, $K\in\{64,128,256\}$ & 8B \\

        Qwen3 8B & 36 & 4096 & TC, Lorsa & $8\times$, $32\times$ &
        TopK, $K=256$ & 8B \\

        \midrule

        Gemma 3 270M & 18 & 640 & Lorsa & 16,384 &
        JReLU, $L_0^\star: 10\!\rightarrow\!20,\; 60\!\rightarrow\!120$ & 8B \\
        
        Gemma 3 1B & 26 & 1152 & Lorsa & 16,384 &
        JReLU, $L_0^\star: 10\!\rightarrow\!20,\; 60\!\rightarrow\!120$ & 8B \\
        
        Gemma 3 4B & 34 & 2560 & Lorsa & 16,384 &
        JReLU, $L_0^\star: 10\!\rightarrow\!20,\; 60\!\rightarrow\!120$ & 8B \\

        \bottomrule
    \end{tabular}
    }
\end{table*}

\begin{table*}[t]
    \caption{
        Lorsa attention configurations used in our workflow scripts.
        \texttt{n\_qk\_heads} is the number of Lorsa query--key heads,
        and $H_{\mathrm{OV}}$ is the number of sparse output-value heads.
        Braced pairs list the values for $8\times$ and $32\times$ dictionaries.
        }
    \centering
    \small
    \resizebox{\textwidth}{!}{%
    \begin{tabular}{lrrrrlll}
        \toprule
        Backbone & Width & \texttt{n\_qk\_heads} & $H_{\mathrm{OV}}$ &
        $d_{\mathrm{QK}}$ & RoPE & Attention & Post-QK norm \\
        \midrule
        Llama 3.2 1B & $8\times,32\times$ & $256,1024$ &
        $16{,}384,65{,}536$ & 64 &
        $\theta=5{\times}10^5$, NTK-by-parts & global & none \\
        Llama 3.2 3B & $8\times,32\times$ & $192,768$ &
        $24{,}576,98{,}304$ & 128 &
        $\theta=5{\times}10^5$, NTK-by-parts & global & none \\
        \midrule
        Qwen3 1.7B & $8\times,32\times$ & $128,512$ &
        $16{,}384,65{,}536$ & 128 &
        $\theta=10^6$ & global & RMS \\
        Qwen3 8B & $8\times,32\times$ & $256,1024$ &
        $32{,}768,131{,}072$ & 128 &
        $\theta=10^6$ & global & RMS \\
        \midrule
        Gemma 3 270M & 16,384 & 64 & 16,384 & 256 &
        $\theta=10^6$ / local $10^4$ & local/global, $w=512$ & RMS \\
        Gemma 3 1B & 16,384 & 64 & 16,384 & 256 &
        $\theta=10^6$ / local $10^4$ & local/global, $w=512$ & RMS \\
        Gemma 3 4B & 16,384 & 64 & 16,384 & 256 &
        $\theta=10^6$ / local $10^4$ & local/global, $w=1024$ & RMS \\
        \bottomrule
    \end{tabular}
    }
    \label{tab:lorsa-head-configs}
\end{table*}

All Lorsa modules with \texttt{initialize\_lorsa\_with\_mhsa} enabled initialize
their QK projections from the corresponding pretrained attention layer. For GQA
backbones, the key projections are expanded from key-value heads to query-head
slots before this copy, so repeated Lorsa K heads share the same initial value
whenever the base model shares a key head. These heads are not tied during
optimization; the shared-K structure is an initialization prior rather than a
hard parameter-sharing constraint.

\subsection{Model Overview}

\subsubsection{Qwen3 Models}

We train per-layer Transcoders and Lorsas for Qwen3-1.7B and Qwen3-8B
\citep{yang2025qwen3}. Both Qwen3 variants use global causal grouped-query
attention (GQA)~\citep{joshua2023gqa} with rotary position embeddings and
per-head QK RMS normalization. We match these details in Lorsa by using the
native QK geometry: Qwen3-1.7B has 16 query heads and Qwen3-8B has 32 query
heads, both with 8 key-value heads, 128-dimensional QK heads, RoPE base $10^6$,
and post-QK RMS normalization. The number of Lorsa QK heads is scaled by the
dictionary expansion factor. Both releases use TopK activation with $8\times$
and $32\times$ dictionary expansion factors. The 1.7B suite includes
$K\in\{64,128,256\}$, while the 8B suite uses $K=256$.

\subsubsection{Llama 3.2 Models}

For Llama 3.2 1B and 3B~\citep{grattafiori2024llama}, we train per-layer
Transcoders and Lorsas across all Transformer layers. Llama 3.2 uses global
causal GQA with rotary position embeddings and Llama-3 long-context rotary
scaling, but does not use QK normalization or local sliding-window attention. We
therefore instantiate Lorsa with the native query-head counts and head
dimensions---32 heads of dimension 64 for the 1B model and 24 heads of
dimension 128 for the 3B model, with 8 key-value heads in both base models---and
disable post-QK normalization. Both model sizes use $8\times$ and $32\times$
dictionaries with $K\in\{64,128,256\}$.

\subsubsection{Gemma 3 Models}

For Gemma 3 270M, 1B, and 4B, we train per-layer Lorsa modules with 16,384
sparse OV features. We use JumpReLU activations and layer-dependent sparsity
targets following Gemma Scope 2~\citep{mcdougall2025gemmascope2}.

Because Gemma uses QK normalization, grouped-query attention, and interleaved
local/global attention~\citep{team2025gemma, joshua2023gqa}, we adapt Lorsa to
match these architectural details, including the corresponding QK head
structure, local attention configuration, and QK normalization. We initialize the 
Lorsa QK parameters from the pretrained Gemma attention weights described above
\citep{wang2026dimensional}, while sharing key projections according to Gemma's GQA structure.

\paragraph{Architecture adaptation for Gemma.}
Gemma attention differs from standard multi-head self-attention in several
important ways: it uses QK normalization, grouped-query attention (GQA), and an
interleaved pattern of local sliding-window and global attention layers, with
different rotary settings across these attention types. We adapt Lorsa to match
these architectural details \citep{team2025gemma, joshua2023gqa}. In
particular, we apply Gemma's QK normalization inside the Lorsa attention
computation, use the appropriate local window and rotary configuration for each
layer, and align the QK head structure with Gemma's GQA layout while keeping
16,384 sparse OV features. The pretrained-attention initialization gives the
sparse decomposition an architecture-matched starting point
\citep{wang2026dimensional}.



\subsection{Activation Data}
For sweeping, training, and analysis (visualization), we use datasets drawn 
from the same distribution to generate activations at the corresponding hook points. 
Llama 3.2 and the Gemma 3 models use 2,048-token  contexts from FineWeb-Edu \citep{penedo2024fineweb}. 
Qwen3 uses 2,048-token contexts from an internal uncropped code--English--Chinese pretraining corpus. 
During training, we apply dataset-wise activation normalization using a fixed scale estimated from the 
training distribution. This brings activations to a stable average norm and improves the numerical scale 
of parameter optimization, while preserving relative token-level magnitudes. After training, we fold the
normalization scale into the learned parameters, allowing the sparse dictionaries to operate directly on 
activations in their original scale at inference time. Cached activations are stored in bfloat16 precision to 
reduce storage and memory overhead. During training, we use mixed precision with bfloat16 autocast
for the forward and loss computation, while sparse-dictionary parameters are stored in float32 unless otherwise 
specified.

\subsection{Training Details}

For the Qwen3 and Llama 3.2 model families, we train TopK Transcoders and
Lorsas. For the Gemma 3 family, Gemma Scope 2 already provides JumpReLU SAEs
and Transcoders; we therefore train only the complementary Lorsas. Below we describe 
the training objectives, initialization and optimization details for the TopK 
Transcoders and Lorsas.

\subsubsection{TopK Transcoders and Lorsas}
\label{sec:topk_transcoder_lorsa}

\paragraph{Objectives} For TopK Transcoders and Lorsas, sparsity is enforced directly by
retaining the $K$ largest positive pre-activations:
\begin{equation}
    \vz
    =
    \operatorname{TopK}\!\left(
        \operatorname{ReLU}(\vh), K
    \right).
\end{equation}
The model is therefore trained using only the reconstruction objective,
\begin{equation}
    \mathcal{L}
    =
    \mathcal{L}_{\mathrm{rec}}
    =
    \mathbb{E}\!\left[
        \left\|
        \vx-\widehat{\vx}
        \right\|_2^2
    \right],
\end{equation}
without additional sparsity penalties.

\paragraph{Initialization}
For TopK Transcoders and Lorsas, we initialize the decoder and encoder weights
from uniform distributions with scale $1/\sqrt{d_{\mathrm{model}}}$ and
$1/\sqrt{d_{\mathrm{model}}m}$, respectively, where $m$ is the dictionary
expansion factor. We do not initialize the encoder as the decoder transpose.
Instead, after sampling the weights, we perform a coarse-to-fine grid search
over the decoder norm on an activation batch and choose the norm that
minimizes the initial reconstruction loss \citep{gao2024topk}. We initialize the decoder bias from
the empirical center of the normalized target activations and initialize the
encoder bias from the mean pre-activation. For Transcoders, we additionally
initialize from the corresponding pretrained MLP layer. For Lorsas, we
initialize the QK parameters from the corresponding pretrained attention layer
and initialize the OV decoder directions within the active subspace of the
training activations \citep{wang2026dimensional}.

\paragraph{Optimization}
For all TopK Transcoders and Lorsas, the initial TopK value is equal to the
final target $K$. Transcoders use SparseAdam, while Lorsas use Adam;
both use $(\beta_1,\beta_2)=(0.9,0.999)$ and no weight decay. We use 5,000 warmup
steps followed by a linear cooldown over the final 20\% of training, with no
gradient clipping.

TopK Transcoders use a batch size of 32,768 tokens. TopK Lorsa batch sizes are
32,768 for Qwen3 and $8\times$ Llama, and 16,384 for $32\times$ Llama. Learning
rates and other hyperparameters are summarized in
Appendix~\ref{app:topk_hyperparameters}.


\subsubsection{JumpReLU Lorsa}
\label{sec:jumprelu_lorsa}

Gemma Scope 2 provides a comprehensive suite of JumpReLU SAEs and Transcoders
for Gemma 3. Following this design, we use JumpReLU activations when training
Lorsas for the Gemma 3 models.

\paragraph{Objectives.}
The reconstruction objective is
\begin{equation}
    \mathcal{L}_{\mathrm{rec}}
    =
    \mathbb{E}\!\left[
        \left\|
        \vx-\widehat{\vx}
        \right\|_2^2
    \right].
\end{equation}
We encourage a target sparsity level $L_0^{*}$ using
\begin{equation}
    \mathcal{L}_{\mathrm{sparse}}
    =
    \lambda
    \frac{2}{L_0^{*}}
    \left(
        L_0-L_0^{*}
    \right)^2 ,
\end{equation}
where $\lambda$ is the sparsity coefficient
\citep{mcdougall2025gemmascope2}.

To mitigate dead features, we additionally use an auxiliary loss
\citep{rajamanoharan2024jumprelu},
\begin{equation}
    \mathcal{L}_{\mathrm{aux}}
    =
    \lambda_{\mathrm{aux}}
    \mathbb{E}\!\left[
        \sum_{i \in \mathcal{D}}
        \max\!\left(0, \tau_i - \tilde{z}_i\right)
    \right],
\end{equation}
where $\mathcal{D}$ denotes the set of dead latents, $\tau_i$ is the detached
JumpReLU threshold, $\tilde{z}_i$ is the corresponding pre-activation, and
$\lambda_{\mathrm{aux}}$ controls the auxiliary loss strength.

The final training objective is
\begin{equation}
    \mathcal{L}
    =
    \mathcal{L}_{\mathrm{rec}}
    +
    \mathcal{L}_{\mathrm{sparse}}
    +
    \mathcal{L}_{\mathrm{aux}} .
\end{equation}

\paragraph{Initialization.}
Gemma Lorsas are initialized from the original multi-head self attention
layers. We copy the query projections directly and initialize the key
projections according to Gemma's GQA structure, so heads sharing a pretrained
key-value head start from the same key projection. The repeated key heads are
then trained independently. We also initialize the OV decoder directions from
the active subspace of the training attention outputs as describe in TopK Lorsa.

\paragraph{Optimization.}
Gemma Lorsas are trained with Adam. We use 5,000 warmup steps and linearly decay 
the learning rate over the final 20\% of training to $1/32$ of its peak value. 
To avoid the excessive number of dead features, we linearly warm up the $\ell_1$ 
coefficient over the first 6\% of training, corresponding to 480M tokens for the 8B model, 
and use an auxiliary loss to further improve dictionary utilization.

For JumpReLU, we use a separate learning rate for the learnable thresholds: $0.1$ for 
settings with $L_0 \in [10, 20]$ and $0.3$ for settings with $L_0 \in [100, 200]$. 
We additionally apply global gradient clipping at $1.0$ to stabilize threshold optimization 
and set the $L_p$ coefficient to $0.003$.

The token batch sizes are 16,384, 32,768, and 131,072 for Gemma 3 270M, 1B,
and 4B, respectively. Learning rates and other hyperparameters are summarized in
Appendix~\ref{app:lorsa-hyperparameters}.

\section{Sparse Decomposition Evaluation}
\label{sec:evaluation}

\subsection{Evaluation Setup and Metrics}

We evaluate the learned sparse decompositions along four dimensions:
dictionary quality, scaling behavior, model-level fidelity, and feature
interpretability.

\paragraph{Dictionary quality.}
We report explained variance (EV), $L_0$, latent density, firing frequency,
and the number of dead latents to characterize reconstruction quality and
feature utilization.

\paragraph{Scaling behavior.}
We evaluate how reconstruction quality changes as the dictionary width and
TopK sparsity level $K$ increase, focusing primarily on the scaling of EV.

\paragraph{Model-level fidelity.}
We insert the trained Transcoders and Lorsas into the model and compare the
replacement model against the original model using output KL divergence and
the change in language-modeling loss,
$\Delta\mathcal{L}_{\mathrm{LM}}$.

\paragraph{Feature interpretability.}
We automatically generate and evaluate feature explanations following prior
automated interpretability protocols
\citep{cunningham2024circuits, karvonen2025saebench}, and report the
corresponding automated interpretability scores.

\paragraph{Feature steering.}
We evaluate the controllability of sparse features through activation 
steering~\citep{rimsky2024steering}. Specifically, we consider two steering tasks: 
(i) text-style transfer, following Qwen-Scope~\citep{deng2026qwenscope}, 
and (ii) personality steering, where we amplify personality-related features 
to induce corresponding behavioral traits~\citep{chen2025persona}.

\subsection{Sparse Dictionary Quality}

We first evaluate the reconstruction--sparsity trade-off of the trained
Transcoders and Lorsas across model families and layers. We report EV and
$L_0$, together with firing-frequency, latent-density, and dead-latent
statistics.

% TODO: Main EV--L0 results across models and layers.
% TODO: Firing-frequency and dead-latent distributions.


\subsection{Scaling Behavior}

We next study how sparse decomposition quality scales with dictionary capacity.
For TopK models, we vary both the dictionary expansion factor and $K$ and
measure the corresponding change in EV. This allows us to examine whether
larger dictionaries and denser operating points consistently improve
reconstruction across model families.

% TODO: EV versus width.
% TODO: EV versus TopK K.
% TODO: Compare scaling trends across Llama 3.2 and Qwen3.


\subsection{Replacement Fidelity}

Reconstruction quality alone does not guarantee that a sparse module faithfully
preserves the original model computation. We therefore replace the original
MLP and attention modules with the corresponding Transcoders and Lorsas and
measure the resulting change in model behavior. We report output KL divergence
and $\Delta\mathcal{L}_{\mathrm{LM}}$, with lower values indicating higher
replacement fidelity.

% TODO: Per-layer replacement results.
% TODO: End-to-end replacement with progressively more layers replaced.
% TODO: KL divergence and delta LM loss.


\subsection{Feature Interpretability}

Finally, we evaluate whether the learned sparse features admit meaningful
natural-language descriptions. We apply automated feature interpretation to
Transcoder and Lorsa features and evaluate the generated explanations using
automated interpretability scores
\citep{cunningham2024circuits, karvonen2025saebench}.

% TODO: Aggregate automated interpretability scores.
% TODO: Representative Transcoder and Lorsa feature examples.



\section{Attribution Graph Construction}
\label{sec:attribution_graph_construction}

We use a per-layer Transcoder to sparsify each MLP sublayer and a Lorsa to
sparsify each attention sublayer. We additionally place a Matryoshka SAE on the
final residual stream as an abstraction-selection interface for graph
construction. Together, these modules provide a sparse replacement of the full
Transformer computation, over which we construct attribution graphs following
circuit tracing~\citep{ameisen2025circuit, shu2026completereplacement}.

Following~\citet{shu2026completereplacement}, we define an attribution graph as
\begin{equation}
    \mathcal{G}
    =
    (\mathcal{V}, \mathcal{E}, \mathcal{A}),
    \label{eq:attribution_graph}
\end{equation}
where the node set consists of
\begin{equation}
    \mathcal{V}
    =
    \mathcal{V}_{\mathrm{embed}}
    \cup
    \mathcal{V}_{\mathrm{transcoder}}
    \cup
    \mathcal{V}_{\mathrm{lorsa}}
    \cup
    \mathcal{V}_{\mathrm{matryoshka}}
    \cup
    \mathcal{V}_{\mathrm{error}}
    \cup
    \mathcal{V}_{\mathrm{logit}} .
    \label{eq:attribution_graph_nodes}
\end{equation}

For a source node $s$ and a downstream target node $t$, their effective
encoder and decoder directions depend on the node type:
\begin{equation}
\begin{gathered}
\vw_{\mathrm{enc},t}
=
\begin{cases}
    \vw_{\mathrm{enc},t},
        & t \text{ is a Transcoder or Matryoshka feature}, \\
    \vw_{V,t},
        & t \text{ is a Lorsa feature}, \\
    \vw_{\mathrm{unembed},t},
        & t \text{ is a logit},
\end{cases}
\\[6pt]
\vw_{\mathrm{dec},s}
=
\begin{cases}
    \vw_{\mathrm{dec},s},
        & s \text{ is a Transcoder or Matryoshka feature}, \\
    \vw_{O,s},
        & s \text{ is a Lorsa feature}, \\
    \vw_{\mathrm{embed},s},
        & s \text{ is a token embedding}, \\
    \epsilon_s,
        & s \text{ is an error node}.
\end{cases}
\end{gathered}
\label{eq:effective_enc_dec}
\end{equation}

The direct attribution from $s$ to $t$ is
\begin{equation}
    \alpha_{s \rightarrow t}
    =
    a_s\,
    P^{t}_{i,j}\,
    \vw_{\mathrm{dec},s}^{\top}
    \vw_{\mathrm{enc},t},
    \label{eq:direct_attribution_score}
\end{equation}
where $a_s$ is the source activation and $P^{t}_{i,j}$ captures the
token-to-token interaction associated with the target computation, with
$P^{t}=I$ for non-attention components.

Starting from a target logit, we recursively trace high-attribution edges
backward through the sparse replacement model to obtain the attribution graph.
In practice, we compute these scores using gradients while detaching attention
patterns and LayerNorm denominators, following~\citet{ameisen2025circuit}.

\subsection{Matryoshka SAEs for Uninterested Circuits}
\label{sec:matryoshka_uninterested_circuits}

Full attribution graphs often contain pathways that are faithful to the target
computation but irrelevant to the analyst's question. We refer to these
components as \emph{uninterested circuits}. Llama Scope 2 addresses this problem
by using a Matryoshka SAE on the final residual stream to choose the abstraction
level before attribution is traced.

Let the Matryoshka widths be
$0=b_0<b_1<\cdots<b_M=d_{\mathrm{SAE}}$. For a residual activation
$\vx_i$ at token position $i$, the Matryoshka SAE produces latent activations
$\vz_i\in\mathbb{R}^{d_{\mathrm{SAE}}}$. During circuit construction, the user
selects a feature range $R=[b_a,b_b)$ whose endpoints are Matryoshka
boundaries. We expose only features in this range by applying the mask
\begin{equation}
    z^{(R)}_{i,k}
    =
    z_{i,k}\,\mathbf{1}\{k\in R\},
    \label{eq:matryoshka_range_mask}
\end{equation}
and use the corresponding residual reconstruction
\begin{equation}
    \widehat{\vx}^{(R)}_i
    =
    \vb_D
    +
    \sum_{k\in R} z_{i,k}\vw_{D,k},
    \qquad
    \epsilon^{(R)}_i
    =
    \vx_i-\widehat{\vx}^{(R)}_i .
    \label{eq:matryoshka_range_reconstruction}
\end{equation}
This induces a range-conditioned attribution graph
$\mathcal{G}_R=(\mathcal{V}_R,\mathcal{E}_R,\mathcal{A}_R)$ whose Matryoshka
nodes are restricted to
\begin{equation}
    \mathcal{V}^{(R)}_{\mathrm{matryoshka}}
    =
    \{(i,k): z_{i,k}\neq 0,\ k\in R\}.
    \label{eq:matryoshka_range_nodes}
\end{equation}
Features outside $R$ are not available as graph nodes; their omitted
contribution remains in the explicit residual reconstruction-error term
$\epsilon^{(R)}_i$. Choosing a prefix $[0,b_j)$ traces the computation at a
cumulative abstraction level, while choosing an adjacent segment
$[b_{j-1},b_j)$ isolates the incremental features added at that level. In this
way, Matryoshka SAEs turn uninterested circuits into a graph-construction
choice rather than a post-hoc pruning problem.

\begin{figure}[t]
    \centering
    \fbox{\rule{0pt}{1.6in}\rule{0.92\linewidth}{0pt}}
    \caption{Overview of Llama Scope 2. The final figure will show how sparse
    autoencoders, transcoders, and Lorsas cover the principal representations
    and computations in a Transformer block.}
    \label{fig:overview}
\end{figure}

\section{Case studies}
\label{sec:case_studies}

\subsection{Factual Recall}

\subsection{Shortcut Reasoning}

\subsection{Model Circuit Comparison}

\subsection{Uninterested Circuit}



\section{Attribution Graph Evaluation}
\label{sec:attribution_graph_evaluation}

We evaluate the quality of attribution graphs using two graph-level metrics,
the \emph{replacement score} and the \emph{completeness score}, following prior
circuit-tracing work~\citep{ameisen2025circuit,shu2026completereplacement}.
Both metrics quantify the extent to which the traced computation is explained
by sparse features rather than reconstruction errors.

\subsection{Graph Scores}

Let $\mathcal{P}_{\mathrm{feat}}$, $\mathcal{P}_{\mathrm{error}}$, and
$\mathcal{P}_{\mathrm{emb}}$ denote attribution paths originating from sparse
feature nodes, reconstruction-error nodes, and input embeddings, respectively,
and let $w(p)$ denote the attribution strength of path $p$.

\paragraph{Replacement score} measures the fraction of all edges from the attribution graph 
that originate from feature nodes rather than error nodes, weighted by their attribution 
strengths:
\begin{equation}
S_{\mathrm{repl}}
=
\frac{
    \sum_{p \in \mathcal{P}_{\mathrm{feat}}} w(p)
}{
    \sum_{p \in \mathcal{P}_{\mathrm{feat}}} w(p)
    +
    \sum_{p \in \mathcal{P}_{\mathrm{error}}} w(p)
}.
\label{eq:replacement_score}
\end{equation}

\paragraph{Completeness score} measures the fraction of total attribution to the 
logit that can be traced all the way back to the input embeddings,
rather than to error nodes:
\begin{equation}
S_{\mathrm{comp}}
=
\frac{
    \sum_{p \in \mathcal{P}_{\mathrm{emb}}} w(p)
}{
    \sum_{p \in \mathcal{P}_{\mathrm{emb}}} w(p)
    +
    \sum_{p \in \mathcal{P}_{\mathrm{error}}} w(p)
}
=
1 -
\frac{
    \sum_{p \in \mathcal{P}_{\mathrm{error}}} w(p)
}{
    \sum_{p \in \mathcal{P}_{\mathrm{emb}}} w(p)
    +
    \sum_{p \in \mathcal{P}_{\mathrm{error}}} w(p)
}.
\label{eq:completeness_score}
\end{equation}

Higher scores indicate that a larger fraction of the traced computation is
accounted for by sparse, interpretable replacement features.
For complete replacement models, these features include both transcoder
features for MLP computation and Lorsa features for attention computation.

\subsection{Perturbation Fidelity}

We evaluate whether the replacement model preserves the downstream effect of
an upstream perturbation. We apply the same perturbation to the original and
replacement models and measure the cosine similarity between their resulting
downstream activation changes. A higher similarity indicates that the
replacement model better preserves the causal propagation of the original model.

\subsection{Attribution--Intervention Consistency}

We further test whether attribution scores predict the actual causal effects of
feature interventions. We steer individual features and measure their effects
on the output logits, then compute the Pearson correlation between these
intervention effects and the corresponding attribution scores. A higher
correlation indicates that the attribution graph more faithfully reflects the
causal influence of its features.


\section{The Llamascopium Framework}
\label{sec:llamascopium}

Llamascopium is the framework used to construct, analyze, and serve the Llama
Scope 2 release. Its design separates sparse-dictionary architectures from data
production, distributed execution, analysis, and visualization, allowing the
same pipeline to support multiple model families and dictionary types.

\subsection{Unified Sparse-Dictionary Interface}

% TODO: Describe SparseDictionaryConfig, SparseDictionary, serialization, and
%       from_pretrained support.
% TODO: List the architectures relevant to this report: SAE, Matryoshka SAE,
%       transcoder, Lorsa, and complete replacement models.

\subsection{Activation Pipeline}

The activation pipeline provides a common interface for generating, loading,
filtering, shuffling, and batching model activations. It supports both flattened
one-dimensional activations and sequence-preserving two-dimensional
activations.
% TODO: Describe ActivationFactory, hook-point configuration, sharding, and
%       model backends.

\subsection{Distributed Runners}

Llamascopium exposes high-level runners for activation generation, training,
evaluation, and feature analysis. These runners support combinations of data,
tensor, model, and head parallelism according to the dictionary architecture.
% TODO: Add a compact mapping from task and dictionary type to the parallel
%       dimensions used in Llama Scope 2.

\subsection{Analysis, Storage, and Visualization}

Feature-level results are stored in a database-backed representation and served
through an analysis API and web interface. Supported analyses include
activation sampling, firing statistics, activation histograms, direct logit
attribution, automated interpretation, and architecture-specific metadata such
as Lorsa source--target patterns.
% TODO: Describe MongoDB/GridFS records and the FastAPI/UI boundary without
%       exposing deployment-specific details.

\subsection{Circuit Analysis Engine}

The circuit engine composes sparse replacement modules with the language model
and traces attribution from output logits through residual, MLP, and attention
latents. Reconstruction-error nodes remain explicit in the resulting graph.
% TODO: Describe greedy attribution collection, OV attribution, Q/K tracing,
%       pruning, and intervention interfaces.

\begin{figure}[t]
    \centering
    \fbox{\rule{0pt}{1.7in}\rule{0.92\linewidth}{0pt}}
    \caption{Planned architecture diagram for Llamascopium, covering model
    backends, activation factories, sparse dictionaries, distributed runners,
    analysis storage, visualization, and circuit tracing.}
    \label{fig:framework}
\end{figure}



\section{Limitations and Open Problems}
\label{sec:limitations}

\paragraph{Dark matter of sparse decomposition.}

\paragraph{Inhibitory Circuits.}

\paragraph{}


% TODO: Add model-family coverage, context-length, and compute limitations.

\section{Conclusion}
\label{sec:conclusion}

Llama Scope 2 extends sparse dictionary releases from representation-level
feature vocabularies toward sparse decompositions of Transformer computation.
This report documents a unified suite for Qwen3, Llama 3.2, and Gemma 3 models
and Llamascopium, the framework used to train, analyze, visualize, and compose
these dictionaries for circuit analysis.
% TODO: Add final checkpoint and code release links.

\section*{Acknowledgments}

% TODO: Add acknowledgments and compute support.

\bibliography{iclr2025_conference}
\bibliographystyle{iclr2025_conference}

\appendix

\section{Released Checkpoints}
\label{app:checkpoints}

% TODO: Add one complete row per released checkpoint family.

\section{Training Hyperparameters}
\label{app:hyperparameters}
\subsection{TopK Training Hyperparameters}
\label{app:topk_hyperparameters}

\begin{table*}[t]
    \centering
    \caption{
    Training hyperparameters for TopK Transcoders and Lorsa modules
    on the Llama 3.2 models.
    }
    \label{tab:topk_training_hyperparameters_llama}

    \captionsetup{skip=10pt}

    \footnotesize
    \renewcommand{\arraystretch}{1.0}
    \setlength{\tabcolsep}{8pt}

    \begin{adjustbox}{max width=\textwidth}
    \begin{tabular}{lllccc}
        \toprule
        Model
        & Width
        & Module
        & $K$
        & Learning rate
        & $\lambda_{1}$ \\
        \midrule

        % ==================== Llama 3.2 1B ====================
        \multirow{12}{*}{Llama 3.2 1B}
        & \multirow{6}{*}{$8\times$}
        & \multirow{3}{*}{Transcoder}
        & 64  & $1{\times}10^{-4}$ & $8{\times}10^{-5}$ \\
        & & & 128 & $1{\times}10^{-4}$ & $8{\times}10^{-5}$ \\
        & & & 256 & $1{\times}10^{-4}$ & $8{\times}10^{-5}$ \\
        \cmidrule(lr){3-6}
        & & \multirow{3}{*}{Lorsa}
        & 64  & $6{\times}10^{-5}$ & $8{\times}10^{-5}$ \\
        & & & 128 & $6{\times}10^{-5}$ & $8{\times}10^{-5}$ \\
        & & & 256 & $7{\times}10^{-5}$ & $8{\times}10^{-5}$ \\
        \cmidrule(lr){2-6}
        & \multirow{6}{*}{$32\times$}
        & \multirow{3}{*}{Transcoder}
        & 64  & $1{\times}10^{-4}$ & $8{\times}10^{-5}$ \\
        & & & 128 & $1{\times}10^{-4}$ & $8{\times}10^{-5}$ \\
        & & & 256 & $1{\times}10^{-4}$ & $8{\times}10^{-5}$ \\
        \cmidrule(lr){3-6}
        & & \multirow{3}{*}{Lorsa}
        & 64  & $2{\times}10^{-4}$ & $8{\times}10^{-5}$ \\
        & & & 128 & $2{\times}10^{-4}$ & $8{\times}10^{-5}$ \\
        & & & 256 & $1{\times}10^{-4}$ & $8{\times}10^{-5}$ \\

        \midrule

        % ==================== Llama 3.2 3B ====================
        \multirow{12}{*}{Llama 3.2 3B}
        & \multirow{6}{*}{$8\times$}
        & \multirow{3}{*}{Transcoder}
        & 64  & $1{\times}10^{-4}$ & $8{\times}10^{-5}$ \\
        & & & 128 & $1{\times}10^{-4}$ & $8{\times}10^{-5}$ \\
        & & & 256 & $9{\times}10^{-5}$ & $8{\times}10^{-5}$ \\
        \cmidrule(lr){3-6}
        & & \multirow{3}{*}{Lorsa}
        & 64  & $6{\times}10^{-5}$ & $8{\times}10^{-5}$ \\
        & & & 128 & $6{\times}10^{-5}$ & $8{\times}10^{-5}$ \\
        & & & 256 & $6{\times}10^{-5}$ & $8{\times}10^{-5}$ \\
        \cmidrule(lr){2-6}
        & \multirow{6}{*}{$32\times$}
        & \multirow{3}{*}{Transcoder}
        & 64  & $1{\times}10^{-4}$ & $8{\times}10^{-5}$ \\
        & & & 128 & $8{\times}10^{-5}$ & $8{\times}10^{-5}$ \\
        & & & 256 & $9{\times}10^{-5}$ & $8{\times}10^{-5}$ \\
        \cmidrule(lr){3-6}
        & & \multirow{3}{*}{Lorsa}
        & 64  & $6{\times}10^{-5}$ & $8{\times}10^{-5}$ \\
        & & & 128 & $6{\times}10^{-5}$ & $8{\times}10^{-5}$ \\
        & & & 256 & $6{\times}10^{-5}$ & $8{\times}10^{-5}$ \\

        \bottomrule
    \end{tabular}
    \end{adjustbox}
\end{table*}

\begin{table*}[t]
    \centering
    \caption{
    Training hyperparameters for TopK Transcoders and Lorsa modules
    on the Qwen3 models. For Qwen3 1.7B,
    $\mathcal{S}_{\mathrm{lr}}=\{0.6,0.7,0.8,0.9,1,2,3,4\}\times10^{-4}$
    denotes the swept learning-rate set.
    }
    \label{tab:topk_training_hyperparameters_qwen}

    \captionsetup{skip=10pt}

    \footnotesize
    \renewcommand{\arraystretch}{1.0}
    \setlength{\tabcolsep}{8pt}

    \begin{adjustbox}{max width=\textwidth}
    \begin{tabular}{lllccc}
        \toprule
        Model
        & Width
        & Module
        & $K$
        & Learning rate
        & $\lambda_{1}$ \\
        \midrule

        % ==================== Qwen3 1.7B ====================
        \multirow{12}{*}{Qwen3 1.7B}
        & \multirow{6}{*}{$8\times$}
        & \multirow{3}{*}{Transcoder}
        & 64  & $\mathcal{S}_{\mathrm{lr}}$ & -- \\
        & & & 128 & $\mathcal{S}_{\mathrm{lr}}$ & -- \\
        & & & 256 & $\mathcal{S}_{\mathrm{lr}}$ & -- \\
        \cmidrule(lr){3-6}
        & & \multirow{3}{*}{Lorsa}
        & 64  & $\mathcal{S}_{\mathrm{lr}}$ & -- \\
        & & & 128 & $\mathcal{S}_{\mathrm{lr}}$ & -- \\
        & & & 256 & $\mathcal{S}_{\mathrm{lr}}$ & -- \\
        \cmidrule(lr){2-6}
        & \multirow{6}{*}{$32\times$}
        & \multirow{3}{*}{Transcoder}
        & 64  & $\mathcal{S}_{\mathrm{lr}}$ & -- \\
        & & & 128 & $\mathcal{S}_{\mathrm{lr}}$ & -- \\
        & & & 256 & $\mathcal{S}_{\mathrm{lr}}$ & -- \\
        \cmidrule(lr){3-6}
        & & \multirow{3}{*}{Lorsa}
        & 64  & $\mathcal{S}_{\mathrm{lr}}$ & -- \\
        & & & 128 & $\mathcal{S}_{\mathrm{lr}}$ & -- \\
        & & & 256 & $\mathcal{S}_{\mathrm{lr}}$ & -- \\

        \midrule

        % ==================== Qwen3 8B ====================
        \multirow{4}{*}{Qwen3 8B}
        & \multirow{2}{*}{$8\times$}
        & Transcoder & 256 & $1{\times}10^{-4}$ & -- \\
        & & Lorsa & 256 & $9{\times}10^{-5}$ & -- \\
        \cmidrule(lr){2-6}
        & \multirow{2}{*}{$32\times$}
        & Transcoder & 256 & $8{\times}10^{-5}$ & -- \\
        & & Lorsa & 256 & $6{\times}10^{-5}$ & -- \\

        \bottomrule
    \end{tabular}
    \end{adjustbox}
\end{table*}

% TODO: Add separate tables for SAEs, transcoders, and Lorsas.

\section{Additional Evaluation Results}
\label{app:additional-results}

% TODO: Add per-layer reconstruction, sparsity, firing-frequency, and
%       replacement metrics.

\section{Analysis and Visualization Details}
\label{app:analysis}

% TODO: Document sampling, histogram construction, DLA, autointerpretation, and
%       circuit graph settings.

\section{Llamascopium API and Reproducibility}
\label{app:llamascopium}

% TODO: Add installation instructions, versioned configuration examples,
%       hardware requirements, and commands for reproducing representative
%       training, analysis, and circuit-tracing runs.

\end{document}
